%! Piecewise \& Step Functions — Unit 2, Lesson 2.4 — Exit Ticket (Answer Key)
\documentclass[10pt]{article}
\usepackage{algebra2-article}
\usepackage{algebra2-key}   % pulls in -boxes; do NOT also load -boxes
\newcommand{\TallMath}[1]{$\displaystyle #1\rule[-1.4em]{0pt}{3.2em}$}
% Coordinate grid: {xmin}{xmax}{ymin}{ymax}
\newcommand{\lgrid}[4]{%
  \draw[step=1, linegray!40, very thin] (#1,#3) grid (#2,#4);
  \draw[->, semithick, charcoal] (#1,0)--(#2,0) node[below right, font=\scriptsize]{$x$};
  \draw[->, semithick, charcoal] (0,#3)--(0,#4) node[left, font=\scriptsize]{$y$};}
% Endpoint dots
\newcommand{\closeddot}[2]{\fill[forest] (#1,#2) circle (3.5pt);}
\newcommand{\opendot}[2]{\draw[very thick, forest, fill=white] (#1,#2) circle (3.5pt);}

\begin{document}

\pageheader{Unit 2, Lesson 2.4}{Exit Ticket --- Answer Key}
\namedateperiod

\vspace{0.10in}

No notes. Show your work in the space provided.

\begin{enumerate}[leftmargin=1.8em, itemsep=14pt]
  \item \textbf{Evaluate.} For $f(x)=\begin{cases} 3x, & x<1 \\ x+4, & x\ge 1 \end{cases}$
        \\[4pt]
        $f(0)={\color{keyred}\mathbf{0}}$ \qquad $f(1)={\color{keyred}\mathbf{5}}$ \qquad
        $f(3)={\color{keyred}\mathbf{7}}$
        \\[4pt]
        Which piece did you use for $f(1)$, and how do you know?
        \ans{the bottom piece $x+4$, because $1\ge 1$ (the boundary belongs to the $\ge$ piece).}

  \item \textbf{Read the graph.} The piecewise function below is shown.
        \begin{center}
        \begin{tikzpicture}[scale=0.5]
          \lgrid{-3}{3}{-2}{4}
          \draw[navy, dashed, semithick] (0,-2)--(0,4);
          \draw[very thick, forest] (-3,-1)--(0,-1);
          \opendot{0}{-1}
          \draw[very thick, forest] (0,1)--(2,3);
          \closeddot{0}{1}
        \end{tikzpicture}
        \end{center}
        $f(-2)={\color{keyred}\mathbf{-1}}$; \quad $f(0)={\color{keyred}\mathbf{1}}$. \quad Is $f$
        continuous at $x=0$? Explain.
        \ans{No --- the graph jumps from $-1$ (open) up to $1$ (closed); the two pieces disagree at $x=0$.}

  \item \textbf{Multiple choice (SOL-style).} The greatest-integer function gives $\lfloor -1.5\rfloor=$
        \begin{enumerate}[label=(\Alph*), leftmargin=1.8em, itemsep=2pt]
          \item $-1$
          \item $-2$ \quad\textcolor{keyred}{\textbf{$\leftarrow$ correct}}
          \item $1$
          \item $2$
        \end{enumerate}
        \begin{itemize}[leftmargin=1.4em, nosep]
          \item Round \emph{down}: the greatest integer $\le -1.5$ is $-2$ (not $-1$, which is larger than $-1.5$).
        \end{itemize}
\end{enumerate}

\begin{teachernote}
Item~1 checks piece selection at a boundary --- students who wrote $f(1)=3$ used the strict-$<$ piece and
missed that $x=1$ belongs to the $\ge$ piece. Item~2 checks reading a graph with a \emph{discontinuity}:
take the \emph{closed} dot for $f(0)=1$, and justify the jump by the pieces disagreeing. Item~3 is the
classic greatest-integer trap: rounding \emph{down} sends $-1.5$ to $-2$. Sort exit tickets into ``got it /
wrong-piece slip / open-vs-closed slip / continuity slip / floor-rounding slip'' to decide whether
Lesson~2.5 needs a two-minute re-teach of piece selection.
\end{teachernote}

\end{document}
